\documentclass[12pt]{article}
\usepackage{amsmath, amssymb, graphicx, hyperref, geometry}
\geometry{a4paper, margin=1in}

\title{\textbf{CNNs from Scratch: Part 2 - Building the Architecture}}
\author{Build in Public Series}
\date{\today}

\begin{document}
\maketitle

\begin{center}
    \includegraphics[width=\textwidth]{Post-2-Banner.png}
\end{center}

\section{Introduction}
In Part 1, we explored the mathematical foundations of convolutions. Now we translate theory into code, building a complete CNN architecture in pure PyTorch.

\section{Building Blocks}

\subsection{Convolutional Layer}
Our custom Conv2D layer implements the sliding window operation:
\begin{itemize}
    \item Learnable kernels with Xavier initialization
    \item Configurable stride, padding, and dilation
    \item Efficient implementation via PyTorch's F.conv2d
\end{itemize}

\subsection{Pooling Layers}
We implement both MaxPool2D and AveragePool2D for spatial downsampling. Max pooling preserves the strongest activations while average pooling smooths the representation.

\subsection{Activation Functions}
ReLU (Rectified Linear Unit) is the default choice:
\begin{equation}
    \text{ReLU}(x) = \max(0, x)
\end{equation}
It introduces non-linearity while avoiding vanishing gradients.

\section{Architecture Design}

\subsection{LeNet-5 (1998)}
One of the earliest successful CNNs, designed for digit recognition:
\begin{verbatim}
Conv1 (6 filters, 5x5) -> AvgPool -> Conv2 (16 filters, 5x5) -> AvgPool
-> FC1 (120) -> FC2 (84) -> Output (10)
\end{verbatim}

\subsection{SimpleCNN}
A modern architecture demonstrating the classic pattern:
\begin{verbatim}
Conv -> ReLU -> Pool -> Conv -> ReLU -> Pool -> FC -> Output
\end{verbatim}

\subsection{DeepCNN}
A VGG-inspired deeper network with reusable CNN blocks that double channels at each stage.

\section{Implementation Details}

\subsection{Parameter Count}
CNNs are parameter-efficient due to weight sharing. A 3x3 convolution with 64 input and output channels has only:
\begin{equation}
    64 \times 64 \times 3 \times 3 = 36,864 \text{ parameters}
\end{equation}
Compare this to a fully connected layer with the same dimensions, which would have millions.

\subsection{Output Size Calculation}
For an input of size $H \times W$:
\begin{equation}
    H_{out} = \left\lfloor\frac{H + 2 \cdot \text{padding} - \text{dilation} \cdot (\text{kernel\_size} - 1) - 1}{\text{stride}} + 1\right\rfloor
\end{equation}

\section{Training Strategy}
We train on MNIST with:
\begin{itemize}
    \item Cross-entropy loss for multi-class classification
    \item Adam optimizer with learning rate scheduling
    \item Dropout regularization in fully connected layers
\end{itemize}

\section{Conclusion}
With our CNN implemented, we're ready to train and visualize what the network learns. Part 3 explores filter visualization and feature map analysis.

\end{document}
